\section*{Course Syllabus}

\subsection*{Instructor and Contact Information}
\begin{itemize}[labelwidth=\widthof{\textbf{Office Hours:}}]
    \item \textbf{Professor:} Dr. Christine Heitsch, Skiles 211B, (404) 894-4758
    \item \textbf{Office Hours:} Mon and Wed, 3:15-3:45 PM
    \item \textbf{Individual Appointments:} To discuss performance or any private concerns, please send a message through Canvas Inbox to make arrangements.
    \item \textbf{Piazza Forum:} For questions regarding course format or content that would benefit all students, please use the Piazza Forum. Postings can be anonymous (to classmates) and should always be respectful and courteous.
    \item \textbf{2106-B Lecture Assistant (LA):} Juli\'an Villaquir\'a. He will hold office hours in the Math Lab (Clough 280). For more information, including the full schedule, refer to \inlinequote{Academic Resources} at \href{https://math.gatech.edu/undergraduate/current-students}{https://math.gatech.edu/undergraduate/current-students}.
\end{itemize}

\subsection*{Course Details}
\begin{itemize}[labelwidth=\widthof{\textbf{Lectures:}}]
    \item \textbf{Lectures:} Mondays and Wednesdays from 2:00-3:15 PM in Skiles 170.
    \item \textbf{Textbook:} G. Chartrand, A. D. Polimeni, and P. Zhang. \textit{Mathematical Proofs: A Transition to Advanced Mathematics}, Fourth Edition, Pearson, 2018. Any prior edition is also acceptable.
    \item \textbf{Course Description:} \inlinequote{An introduction to proofs in advanced mathematics, intended as a transition to upper division courses including Abstract Algebra I and Analysis I.}
    \item \textbf{Prerequisites:} Math 1552: Integral Calculus or equivalent.
    \item \textbf{Course Topics:}
    \begin{itemize}
        \item Fundamentals of mathematical abstraction: sets (Ch. 1), logic (Ch. 2), equivalence relations (Ch. 9), and functions (Ch. 10).
        \item Thorough development of basic proof techniques (Ch. 3-8): direct, contrapositive, existence, contradiction, and induction.
        \item Introduction to proofs in algebra (Ch. 15) and analysis (Ch. 14).
        \item Selected additional topics (from Ch. 0, 11, 12 as well as partial orders, upper/lower bounds, and $\limsup / \liminf$).
    \end{itemize}
    \item \textbf{Learning Outcomes:} Students should be able to understand and prove simple mathematical statements. For instance, in analysis, students will be able to prove that a given sequence diverges to infinity. In algebra, students will be able to prove that up to isomorphism there are two distinct groups of order four.
\end{itemize}

\subsection*{Grading}
\begin{itemize}[labelwidth=\widthof{\textbf{Grading Scheme:}}]
    \item \textbf{Grading Scheme:} Grades will be based on exams (80\%) and assignments (20\%) according to the standard scale: A (90-100), B (80-89), C (70-79), D (60-69), F (0-59).
\end{itemize}

\subsubsection*{Exams}
\begin{itemize}
    \item There will be \textbf{four hourly exams}, each worth \textbf{20\% of the final grade}.
    \item Exams will be closed book, closed notes, no calculator, and individual tests. No makeup exams are scheduled.
    \item Three exams will be in-class. Exam dates will be reconfirmed at least a week in advance.
    \item \textbf{Hourly 1:} Scheduled for Wednesday, September 10.
    \item \textbf{Hourly 2:} Scheduled for Wednesday, October 8.
    \item \textbf{Hourly 3:} Scheduled for Wednesday, November 5.
    \item \textbf{Hourly 4:} Will be during the final exam period from 2:40-3:55 PM on Friday, December 5.
\end{itemize}

\subsubsection*{Assignments}
\begin{itemize}
    \item There will be \textbf{three types of assignments} -- in-class, online, and homework -- worth a total of \textbf{20\% of the final grade}.
    \item Assignments will be due regularly, including on the \inlinequote{Final Instructional Class Day}.
    \item To succeed in this class, you must be able to demonstrate mastery of the material; assignments are a chance to practice for the exams.
\end{itemize}

\paragraph*{ICAs (In-Class Assignments)}
\begin{itemize}
    \item Given most days, worth \textbf{5\% of the total course grade}.
    \item Graded on a scale of 0-2 based on effort only.
    \item The lowest 2 scores will be dropped when computing the average.
\end{itemize}

\paragraph*{WeBWorK (Online Problem Sets)}
\begin{itemize}
    \item Assigned through Canvas, worth \textbf{5\% of the total course grade}.
    \item No partial credit will be given.
    \item The lowest score will be dropped when computing the average.
\end{itemize}

\paragraph*{Homework (Written Assignments)}
\begin{itemize}
    \item Assigned most weeks, worth \textbf{10\% of the total course grade}.
    \item Typically due one week later.
    \item Solutions must be turned in via Gradescope in Canvas.
    \item Late homework will be accepted with a penalty of 5\% per 6 hours, up to a total of 36 hours.
    \item If 85\% of students complete the class CIOS survey, one homework score will be dropped from the grading scheme computation; otherwise, all scores will be used.
    \item A subset of the assigned problems will be selected for grading. For full credit, all directions must be followed. Illegible and/or unintelligible solutions will receive no credit.
    \item \textbf{Collaboration with classmates is allowed} (and strongly encouraged!) when working on homework problems and other coursework.
    \item \textbf{Content Guidelines for Solutions:}
    \begin{enumerate}
        \item Be written independently in a student's own words.
        \item Clearly acknowledge any person with whom any part of the assignment was discussed.
        \item Properly credit ANY\footnote{GenAI, e.g. ChatGPT, MathGPT, etc., is most definitely an \inlinequote{outside resource}.} outside resource consulted in completing the assignment.
    \end{enumerate}
    \item Any solution that violates these content guidelines will receive no credit. Flagrant or repeated violations will be dealt with as a matter of academic integrity.
    \item \textbf{\LaTeX\ Bonus:} Homework solutions typeset in \LaTeX\ will receive an additional 10\% bonus points, up to a maximum of 100\% per assignment. Problems will be provided in a *.tex source file. Georgia Tech provides access to Overleaf Professional, an online browser-based \LaTeX\ editor: \href{https://www.overleaf.com/edu/gatech}{https://www.overleaf.com/edu/gatech}.
\end{itemize}

\subsection*{Policies}
\begin{itemize}[labelwidth=\widthof{\textbf{Academic Integrity:}}]
    \item \textbf{Attendance:} Regular participation is expected. Exceptions will be accommodated only for valid, documented reasons including (1) official representation of the Institute and (2) medical or personal emergencies (as verified by the Dean of Students).
    \item \textbf{Exceptions:} Any student who may not be able to meet the requirements of the class as stated must contact the professor individually within the first two weeks of class.
    \item \textbf{Academic Integrity:} Students are expected to abide by the GT Academic Honor Code, Student Code of Conduct, and other policies available online through the Office of Student Integrity. Any violations must be reported directly to the Dean of Students.
    \item \textbf{Updates:} This syllabus is subject to modification. Any changes will be announced in Canvas.
\end{itemize}